\section{Introduction}
\label{sec:introduction}

Autonomous agents powered by Large Language Models (LLMs) are increasingly deployed to execute multi-turn, long-horizon tasks, including software development~\cite{wu2025gcc}, multi-session customer support~\cite{deshmukh2024mem0}, scientific exploration, and interactive simulations~\cite{park2023generative}. As these agents interact with users and external environments over days or weeks, they accumulate extensive conversational histories, tool executions, and evolving operational constraints. Managing this accumulated state poses a fundamental challenge: while the agent must retain knowledge of foundational decisions made in early sessions, it must also adapt when those decisions are superseded by newer instructions, all while remaining strictly bounded by finite context budgets~\cite{packer2023memgpt,liu2024lost}.

Historically, agent state management has relied on a continuum of heuristic approximations. Simple \textit{sliding windows} preserve immediate conversational flow but suffer catastrophic forgetting of early project constraints. \textit{Running summarization} compresses conversational history into natural language preambles, but inevitably suffers from information flattening, semantic drift, and irreversible loss of exact technical specifications. \textit{Retrieval-Augmented Generation} (RAG)~\cite{lewis2020retrieval,guu2020realm} and vector-based memory layers~\cite{deshmukh2024mem0} index historical snippets into dense embedding spaces; however, standard retrieval models treat all chunks as timeless truths. When a user changes an instruction (e.g., migrating from PostgreSQL to MongoDB), vector similarity search frequently retrieves both the outdated initial configuration and the updated requirement, introducing conflicting ``ghost context'' that degrades model reasoning~\cite{stale2025benchmark,wu2024longmemeval}.

Recent architectural advances have significantly expanded the nominal context windows of modern foundation models, extending capacity from 4,096 tokens to 128k, 1M, and 2M tokens~\cite{brown2020language,dong2023bamboo}. However, empirical research demonstrates that larger context windows do not solve the memory challenge~\cite{liu2024lost,maharana2024locomo}. First, stuffing massive histories into prompts incurs severe linear or quadratic inference latency and high financial costs. Second, frontier models exhibit the ``Lost-in-the-Middle'' phenomenon, where factual recall deteriorates when critical information is placed in the center of long contexts~\cite{liu2024lost}. Third, in-context prompt stuffing provides no persistence across independent inference calls and offers zero mechanisms for historical auditing or branch exploration. In short, \textbf{large context capacity does not eliminate the necessity for principled context curation}.

Most critically, existing persistent memory architectures commonly employ \textit{destructive update mechanics}~\cite{packer2023memgpt,deshmukh2024mem0,zhou2025ltmsecurity}. When new facts arrive, frameworks typically execute in-place text replacements or database \texttt{UPDATE} operations, permanently erasing previous configurations. When an agent produces an unexpected output or hallucinates, human auditors have no mechanism to reconstruct what the agent believed at that specific point in time. Destructive updates also render agents incapable of answering retrospective historical queries (e.g., ``What was our database target before Tuesday's meeting?'') after factual supersession has occurred.

To address these compounding limitations, we propose \textbf{ManiGit}, an open architecture providing transparent, versioned context memory for long-running AI agents. ManiGit decouples persistent memory accumulation from bounded prompt context assembly by integrating:
\begin{enumerate}
    \item \textbf{An Immutable Message Archive:} An append-only SQLite repository that permanently preserves raw conversation turns with content hashing and session provenance.
    \item \textbf{A Directed Acyclic Graph (DAG) Commit Substrate:} A Git-inspired commit tree where each factual state change is recorded as an immutable, parent-linked commit object, enabling complete auditability and non-destructive historical checkout.
    \item \textbf{Hybrid Rule-First and Confidence-Calibrated Memory Intake:} An intake pipeline that routes unambiguous technical constraints through zero-latency pattern rules, while passing ambiguous candidate assertions to an LLM judge calibrated with explicit confidence thresholds.
    \item \textbf{A Four-Tier Memory Lifecycle Model:} Classifying facts into functional lifecycle states (\texttt{VERBATIM}, \texttt{SUMMARY}, \texttt{ARCHIVED}, \texttt{TOMBSTONE}), isolating superseded or refuted facts from active prompt compilation while retaining them for historical reconstruction.
    \item \textbf{Automated Contradiction Supersession and Why Reports:} Automatically detecting when a new fact contradicts an active fact, archiving the superseded fact, and generating an auditable Why Report recording the classification rule, confidence score, and supersession target.
    \item \textbf{Token-Budget Context Compilation:} Formulating prompt context assembly as a constrained knapsack problem, solved via a greedy score-density compiler that enforces strict compliance with arbitrary token budgets $B$.
\end{enumerate}

Through a rigorous controlled evaluation comprising 880 experimental trials across 11 scenarios, 8 methods/variants, and 5 token budgets ($B \in [128, 2048]$), we demonstrate that ManiGit achieves 86.4\% fact recall [95\% CI: 78.7\%, 91.6\%] while consuming 43.0\% fewer context tokens than Full History and consistently maintaining 100\% compliance with token budgets. Furthermore, in our evaluated benchmark scenarios, ManiGit achieves 100\% accuracy on point-in-time historical reconstruction queries where baseline RAG and sliding windows score 0\%. Controlled ablations quantitatively isolate the contribution of contradiction supersession (-5.0\% degradation when removed), confidence thresholding (-5.0\% degradation when removed), and tier-aware ranking (-1.0\% overall degradation, dropping to 0\% recall on targeted tier competition at $B=128$). The evaluation pipeline, experiment configurations, and raw benchmark artifacts are organized to support reproducibility.
